
\begin{proposition}[Maximum components without a minimality assumption]\label{cdp:st:C.2.UN.69}
Let $f:U\to\Delta$ be a proper degeneration of connected complex
two-tori from a smooth threefold. Let its full central fibre
$S=\sum_vT_v$ be reduced and simple normal crossing, without triple
points. No minimality assumption is made on any $T_v$.
Write $k_v$ for the actual canonical evaluation orders.

At a local maximum $T_v$ of these orders, precisely one of the
following descriptions applies:
\begin{enumerate}
\item Every adjacent component has order $k_v$.
\item The component is an unblown terminal ruled surface meeting
its parent in a section, with order drop one, or an unblown
terminal plane meeting its parent in a line, with order drop two.
These are the actual contractions of \ref{cdp:st:C.2.UN.64}.
\item The component is an unblown terminal Hopf surface, with
the exact axis normal order and positive canonical drop of
\ref{cdp:st:C.2.UN.50} and \ref{cdp:st:C.2.UN.58}.
\end{enumerate}
In the second alternative the ruled base has positive genus.
Thus, when the orders are not all equal, a global maximum is
one of the two kinds of strict-drop leaves.

Moreover every rational double curve can be removed by a finite
sequence of the actual plane-line contractions of \ref{cdp:st:C.2.UN.64}.
At each step the other double curves and their neighbourhoods
are unchanged. If $r_0$ is the number of connected rational
double curves of the original support, the resulting actual
torus degeneration has no rational double curve and satisfies
\[
             e(S)=e(S_{\rm nr})+2r_0 .
\tag{C.NM.1}\label{cdp:eq:C.NM.1}
\]
There are exactly $r_0$ point contractions. Reversing their
recorded ordinary blowups recovers the original support,
including its original smooth threefold neighbourhood.
In particular the rational double edges are removed as terminal
edges of the actual incidence graph; none belongs to a cycle.
\end{proposition}

\begin{proof}
\par\medskip\noindent\textbf{The exact restrictions at an edge.}
Canonical evaluation, adjunction and the actual equation $xy=t$
give the same line maps as (C.MR.3)--(C.MR.4), without a
minimality assumption:
\[
 \begin{split}
 K_{T_v}&=\operatorname{div}(\omega_v)
       =\sum_{e=(v,w)}(k_w-k_v-1)C_e,\\
 N_{C_e/T_v}\otimes N_{C_e/T_w}&\simeq\OO_{C_e},\\
 K_{C_e}&\simeq N_{C_e/T_v}^{\,k_w-k_v}.
 \end{split}
\tag{C.NM.2}\label{cdp:eq:C.NM.2}
\]
The last equality uses that the other boundaries are disjoint
from $C_e$. In particular, for $g_e=g(C_e)$,
\[
                  (k_w-k_v)C_e^2=2g_e-2.
\tag{C.NM.3}\label{cdp:eq:C.NM.3}
\]
These formulas retain the actual meromorphic section and
dual-normal trivialization.

\par\medskip\noindent\textbf{A positive-square rational curve on a component.}
We shall use the following implication for any component, not
only a maximum:
\[
 \text{a smooth rational curve of positive square on }T
                         \quad\Longrightarrow\quad T
                         \text{ is rational}.
\tag{C.NM.4}\label{cdp:eq:C.NM.4}
\]
Here is the comparison with its actual minimal-model map.
For a nonrational ruled minimal model, every rational curve
maps to a point of the positive-genus ruling base. A surviving
such curve is a ruling fibre, of square zero; its strict
transform has nonpositive square. Every exceptional curve
also has negative square. On the Hopf minimal model supplied
by \ref{cdp:st:C.2.UN.19} there are no rational curves, since a map from
$\PP^1$ lifts to $\C^2\setminus\{0\}$ and has constant
coordinate functions. Hence a rational curve on a blowup of
that model is exceptional and has negative square.

The remaining actual minimal models have Kodaira dimension
zero or one by \ref{cdp:st:C.2.UN.15}. In this multiplicity-one setting \ref{cdp:st:C.2.UN.31}
gives a torus or the smooth-multiple elliptic model. Its
canonical line is numerically nonnegative on every curve:
on a torus it is trivial; in the elliptic case its positive
power is the pullback of the actual base line of degree
\[
                   2g(B)-2+\sum_b(1-1/m_b)\ge0.
\tag{C.NM.5}\label{cdp:eq:C.NM.5}
\]
The inequality follows from the pluricanonical pushdown and
the nonnegative Kodaira dimension. For its point-blowup map
$\pi:T\to T_{\min}$, one has
$K_T=\pi^*K_{T_{\min}}+\sum b_EE$, with $b_E>0$.
A nonexceptional rational curve $C$ would therefore have
$K_T\cdot C\ge0$. Positive square instead gives
$K_T\cdot C=-2-C^2<0$ by adjunction. An exceptional curve
has negative square. This proves \eqref{cdp:eq:C.NM.4} through the
actual minimal-model morphism.

If $C\subset T$ is rational with square $s=1$ or $2$,
blow up $s$ distinct smooth points of $C$. The strict
transform $C'$ has square zero. Since $T$ is rational,
$H^1(\OO)=0$ on this blown-up surface. The restriction
sequence for $\OO(C')$, with
$\OO_{C'}(C')\simeq\OO_{\PP^1}$, gives a second section
nonzero everywhere on $C'$. Together with the canonical
section cutting out $C'$ it gives an actual basepoint-free
pencil
\[
                         q:T'\longrightarrow\PP^1.
\tag{C.NM.6}\label{cdp:eq:C.NM.6}
\]
There are no basepoints outside $C'$ because the canonical
section is nonzero there. The fibre $C'$ is reduced and
connected. In the Stein factorization a finite map of
degree greater than one would either give more than one
point above this value or a multiplicity greater than one
on $C'$. Thus $q$ has connected fibres. Adjunction gives
genus zero on a smooth fibre. The ordinary relative
minimal-model contractions identify this rational fibration
with a projective-line bundle followed by point blowups
in fibres; in particular every irreducible vertical curve
is rational. Every original curve disjoint from $C$
is untouched by the selected blowups and has intersection
zero with $q^*\OO_{\PP^1}(1)=\OO(C')$, so it is vertical.
This supplies the precise pencil used below.

\par\medskip\noindent\textbf{A maximum with a rational boundary is already minimal.}
At a local maximum put $a_e=k_v-k_w+1\ge1$. Then
\[
             -K_T=\sum_e a_eC_e,\qquad
             (1-a_e)C_e^2=2g(C_e)-2.
\tag{C.NM.7}\label{cdp:eq:C.NM.7}
\]
The boundary is nonempty in the reducible case. A nonzero
pluricanonical section multiplied by the corresponding
power of this anticanonical section would be a nonzero
holomorphic function having zeros on a compact connected
surface. Consequently $\kappa(T)=-\infty$.

For a rational boundary, \eqref{cdp:eq:C.NM.7} gives exactly
\[
                     (a_e,C_e^2)=(3,1)\ \text{or}\ (2,2).
\tag{C.NM.8}\label{cdp:eq:C.NM.8}
\]
Equation \eqref{cdp:eq:C.NM.4} first makes $T$ rational; no rational
pencil is used before this conclusion. The actual pencil
\eqref{cdp:eq:C.NM.6} then makes every other disjoint boundary rational.
Every such boundary would again have positive square by
\eqref{cdp:eq:C.NM.7}, whereas the Hodge index theorem permits at most
one positive-square class among mutually disjoint curves.
Thus there is only one boundary $C$, and $K_T=-aC$.
Any smooth rational $(-1)$-curve $E$ would satisfy
$a(C\cdot E)=1$, which is impossible for $a=2$ or $3$.
The surface is therefore minimal. The terminal-pair
classification and actual maps in \ref{cdp:st:C.2.UN.64} give its ruled
section or plane-line contraction.

\par\medskip\noindent\textbf{All other blowdowns occur only at coefficient-one elliptic boundaries.}
Suppose now that no boundary is rational. A $(-1)$-curve
$E$ is not in the boundary, and adjunction gives
\[
                     1=(-K_T)\cdot E
                              =\sum_e a_e(C_e\cdot E).
\tag{C.NM.9}\label{cdp:eq:C.NM.9}
\]
Every summand is a nonnegative integer. Exactly one is
one: $E$ meets one boundary transversely once, that
boundary has coefficient one, and $E$ misses all others.
The coefficient-one boundary is elliptic by \eqref{cdp:eq:C.NM.7}.
Its ordinary surface blowdown keeps this boundary smooth
and keeps all boundary curves disjoint. Pushforward of
the actual meromorphic canonical section gives the same
canonical divisor with the same coefficients on the
new surface. This is a statement about the surface
morphism; no ambient contraction of $E$ is invoked.
Repeat to the minimal model.

The minimal possibilities with \eqref{cdp:eq:C.NM.7} are exactly those
computed in (C.MR.6)--(C.MR.7) and \ref{cdp:st:C.2.UN.64}. On a ruled
surface, intersection with a ruling fibre gives the three
partitions of two: one coefficient-two section, one
coefficient-one bisection, or two coefficient-one sections.
The connected elliptic bisection over an elliptic base is
excluded by the actual residue/trace map (C.EL.7)--(C.EL.8).
On the plane the possibilities are a coefficient-three
line or a coefficient-one cubic. On an elliptic Hopf
surface the invariant relative differential gives
$\sum (k_w-k_T)/m_e=\operatorname{val}(T)-2$ with every
summand nonpositive; thus it is either a strict-drop
one-boundary leaf or a two-boundary equal-height component.
The nonlinear Hopf axis formula gives the same strict
leaf alternative, including its exact secondary framing.
For an algebraic-dimension-zero diagonal Hopf surface,
the two-axis character equation and $|\alpha|,|\beta|<1$
make both coefficients one; a single axis would impose
an integral resonance. These are actual marked pairs
by \ref{cdp:st:C.2.UN.50}, \ref{cdp:st:C.2.UN.58} and the proof of \ref{cdp:st:C.2.UN.66}.

In particular whenever a coefficient is greater than
one, there is just one boundary, with that coefficient.
No coefficient-one boundary then exists. Since the
surface blowdowns above preserve every boundary and
its coefficient, \eqref{cdp:eq:C.NM.9} shows that no such blowdown
could have taken place. The original strict-drop maximum
was already minimal. In all remaining cases every
coefficient is one and every neighbour has the same
height. A component at the global maximum with this
property cannot be adjacent to a strict-drop maximum
leaf. Its equal-height neighbours have the same property,
so connectedness propagates it to the whole support.
This proves the maximum assertion, except for the
rational ruled-section exclusion proved next.

\par\medskip\noindent\textbf{A rational section centre cannot lie in the smooth torus-fibre locus.}
In a ruled-section contraction of \ref{cdp:st:C.2.UN.64} with base
$B\simeq\PP^1$, the actual descended family is a
submersion near the compact centre $B$. Its containing
central component $P'$ is smooth there and
\[
                    N_{B/P'}\simeq K_B\simeq\OO_{\PP^1}(-2).
\tag{C.NM.10}\label{cdp:eq:C.NM.10}
\]
Choose a relatively compact tubular neighbourhood of
$B$ inside that smooth central surface and a slightly
larger compact neighbourhood in the submersion locus.
A smooth metric gives horizontal lifts of the two
coordinate vector fields on the disc. For sufficiently
small time their flows exist on the whole smaller
neighbourhood and embed it in an actual nearby fibre.
The induced diffeomorphism preserves the oriented
normal bundle and its Euler number. Thus the nearby
two-torus would contain an embedded sphere of
self-intersection minus two. Its universal cover is
the contractible space $\C^2$, so this sphere map is
nullhomotopic, its homology class is zero, and its
self-intersection is zero. This contradiction excludes
the rational ruled-section alternative.

\par\medskip\noindent\textbf{Following rational edges gives actual point contractions.}
Take a rational double curve and orient it toward its
higher-height endpoint. Equation \eqref{cdp:eq:C.NM.3} makes the
height difference nonzero; at that higher endpoint
its square is positive, equal to one or two. By
\eqref{cdp:eq:C.NM.4} the endpoint surface is rational. The pencil
\eqref{cdp:eq:C.NM.6} makes every other boundary on it rational.
Hodge index leaves at most one positive-square
boundary, the incoming one. Every other boundary
therefore has negative square by \eqref{cdp:eq:C.NM.3}, and its
neighbour has strictly higher height. Follow one
of those edges when it exists.

The heights strictly increase and there are finitely
many components. The path ends at a rational
component having exactly its one incoming boundary.
It has $K=-aC$, with $(a,C^2)$ as in \eqref{cdp:eq:C.NM.8}.
The same integer-intersection argument excludes every
$(-1)$-curve, so it is an actual \ref{cdp:st:C.2.UN.64} leaf. Its ruled
section possibility has just been excluded by
\eqref{cdp:eq:C.NM.10}. It is therefore a plane-line leaf, and
\ref{cdp:st:C.2.UN.64} contracts it to a smooth point of the smooth
central locus of its parent.

Every other boundary misses that line, so the actual
blowdown is an isomorphism on its neighbourhood.
It removes exactly one rational double curve, creates
no new double curve or triple point, retains the
punctured torus family, and decreases $e(S)$ by two.
Apply the argument to the new actual degeneration
until no rational double curve remains. Exactly
$r_0$ steps have been made; \eqref{cdp:eq:C.NM.1} and the recorded
ordinary-blowup reconstruction follow.
\end{proof}


